\section{Taylor and Maclaurin Series}

\begin{definition}[Taylor series]
If $f$ has derivatives of all orders at $a$, its Taylor series at $a$ is
\[ T_f(x; a) = \sum_{n=0}^\infty \frac{f^{(n)}(a)}{n!} (x - a)^n. \]
A \emph{Maclaurin series} is the special case $a = 0$.
\end{definition}

\subsection{Taylor's theorem with remainder}
For $f$ that is $(n+1)$-times differentiable on an interval containing $a$ and $x$:
\[
f(x) = \sum_{k=0}^n \frac{f^{(k)}(a)}{k!} (x - a)^k + R_n(x),
\]
with the Lagrange form of the remainder
\[ R_n(x) = \frac{f^{(n+1)}(\xi)}{(n+1)!} (x - a)^{n+1}, \quad \xi \text{ between } a \text{ and } x. \]
If $R_n(x) \to 0$ as $n \to \infty$, then $f$ equals its Taylor series on that interval.

\subsection{Essential Maclaurin series}
\begin{keybox}[Memorize --- all centered at $0$]
\[
\renewcommand{\arraystretch}{1.6}
\begin{array}{l l l}
e^x = \displaystyle\sum_{n=0}^\infty \dfrac{x^n}{n!} & \text{all } x \\
\sin x = \displaystyle\sum_{n=0}^\infty \dfrac{(-1)^n x^{2n+1}}{(2n+1)!} & \text{all } x \\
\cos x = \displaystyle\sum_{n=0}^\infty \dfrac{(-1)^n x^{2n}}{(2n)!} & \text{all } x \\
\dfrac{1}{1 - x} = \displaystyle\sum_{n=0}^\infty x^n & |x| < 1 \\
\ln(1 + x) = \displaystyle\sum_{n=1}^\infty \dfrac{(-1)^{n+1} x^n}{n} & -1 < x \le 1 \\
\arctan x = \displaystyle\sum_{n=0}^\infty \dfrac{(-1)^n x^{2n+1}}{2n+1} & |x| \le 1 \\
(1 + x)^\alpha = \displaystyle\sum_{n=0}^\infty \binom{\alpha}{n} x^n & |x| < 1
\end{array}
\]
where $\binom{\alpha}{n} = \tfrac{\alpha (\alpha - 1)\cdots(\alpha - n + 1)}{n!}$ (binomial series).
\end{keybox}

\subsection{Using known series}
Almost every Taylor problem on an exam is solved by manipulating the table above:
substitute, differentiate, integrate, or multiply.

\begin{example}[Substitution]
$\displaystyle e^{-x^2} = \sum_{n=0}^\infty \frac{(-1)^n x^{2n}}{n!}.$
\end{example}

\begin{example}[Indefinite integration via series]
$\displaystyle \int e^{-x^2} \dx = C + \sum_{n=0}^\infty \frac{(-1)^n x^{2n+1}}{n! (2n+1)}.$
\end{example}

\begin{example}[Limits via series]
$\displaystyle \lim_{x \to 0} \frac{\sin x - x}{x^3}$:
$\sin x = x - \tfrac{x^3}{6} + O(x^5)$, so numerator $= -x^3/6 + O(x^5)$, ratio $\to -1/6$.
\end{example}

\subsection{Error estimation}
For an alternating Taylor series with terms decreasing in magnitude, the error after $N$ terms
is at most the first omitted term (alternating series remainder).
Otherwise use the Lagrange remainder bound:
\[ |R_n(x)| \le \frac{M}{(n+1)!} |x - a|^{n+1}, \qquad M = \max_{|\xi - a| \le |x - a|} |f^{(n+1)}(\xi)|. \]

\begin{pitfallbox}
A function can have a Taylor series at $a$ that converges everywhere but does \emph{not} equal
the function (e.g.\ $f(x) = e^{-1/x^2}$ at $0$). Equality requires $R_n(x) \to 0$.
\end{pitfallbox}
